\documentclass[10pt,letterpaper]{article}

\usepackage[T1]{fontenc}
\usepackage[utf8]{inputenc}
\usepackage{lmodern}
\usepackage[margin=0.82in,headheight=14pt]{geometry}
\usepackage{microtype,amsmath,amssymb,array,tabularx,longtable,booktabs,multicol}
\usepackage{graphicx,xcolor,enumitem,fancyhdr,hyperref,titlesec,needspace}
\definecolor{Navy}{HTML}{122A42}
\definecolor{Ocean}{HTML}{147D91}
\definecolor{Soft}{HTML}{E9F2F5}
\definecolor{Ink}{HTML}{263745}
\definecolor{Gray}{HTML}{5F6873}
\hypersetup{colorlinks=true,linkcolor=Ocean,urlcolor=Ocean,bookmarksnumbered=true}
\setlength{\parindent}{0pt}
\emergencystretch=2em
\setlength{\parskip}{5.3pt}
\setlist{topsep=3pt,itemsep=2pt,parsep=0pt,leftmargin=*}
\titleformat{\section}{\large\bfseries\color{Navy}}{\thesection}{0.65em}{}
\titleformat{\subsection}{\normalsize\bfseries\color{Ocean}}{\thesubsection}{0.6em}{}
\titleformat{\subsubsection}{\normalsize\bfseries\color{Navy}}{\thesubsubsection}{0.5em}{}
\pagestyle{fancy}\fancyhf{}
\fancyhead[L]{\textsf{Architecture Documentation \& Engineering}}
\fancyhead[R]{\textsf{\DocCategory}}
\fancyfoot[L]{\footnotesize\textcolor{Gray}{Controlled learning / engineering artifact}}
\fancyfoot[R]{\footnotesize\thepage}
\renewcommand{\headrulewidth}{0.35pt}
\newcommand{\DocCategory}{Suite}
\newcommand{\ArchTitle}[4]{%
\renewcommand{\DocCategory}{#3}%
{\small\sffamily\bfseries\color{Ocean} #3 \hfill #4}\par\vspace{7pt}
{\Large\bfseries\color{Navy}#1}\par\vspace{4pt}
{\normalsize\color{Gray}#2}\par\vspace{11pt}
\textcolor{Ocean}{\rule{\linewidth}{1.15pt}}\vspace{4pt}
}
\newcommand{\Note}[1]{\par\smallskip\noindent\fcolorbox{Ocean}{Soft}{\parbox{0.94\linewidth}{\small\textbf{Engineering note.}~#1}}\par\smallskip}
\newcommand{\Outcome}[1]{\textbf{Outcome.} #1\par}
\newcommand{\Evidence}[1]{\textbf{Required evidence.} #1\par}
\newcommand{\BlankLines}[1]{\par\begingroup\setlength{\parskip}{9pt}\count255=0\loop\ifnum\count255<#1\noindent\rule{\linewidth}{0.35pt}\par\advance\count255 by 1\repeat\endgroup\vspace{4pt}}
\newcommand{\BlankBox}[1]{\noindent\fbox{\parbox[t][#1][t]{0.98\linewidth}{\mbox{}}}\par\smallskip}
\newcommand{\Field}[2]{\Needspace{3\baselineskip}\textbf{#1}\par\textit{\color{Gray}#2}\par\BlankLines{2}}
\newcommand{\FieldBlank}[1]{\Needspace{3\baselineskip}\textbf{#1}\par\BlankLines{2}}
\newcommand{\DocProvenance}{\textit{Source foundation: Clements et al., \emph{Documenting Software Architectures: Views and Beyond}, 2nd ed. (2011). Interpretive examples, exercises, checklists, and templates are original adaptations, not reproductions.}\par}
\newcommand{\SmallHeader}[1]{\smallskip\textbf{\color{Navy}#1}\par}
\newcolumntype{Y}{>{\raggedright\arraybackslash}X}
\newcommand{\Tick}{\(\square\)\,}

\begin{document}
\ArchTitle{Software Architecture Documentation: Engineering Reference}{Structures, viewpoints, views, rationale, correspondence, review, and controlled evolution}{Technical Reference}{ARCH-REF-001}
\DocProvenance
\section{Purpose, boundaries, and intellectual provenance}
This manual provides a repeatable method to select, construct, connect, evaluate, publish, and maintain software architecture documentation. It preserves the central Views and Beyond distinction between views of designed structures and additional information spanning multiple views. Examples using infrastructure as code, cloud resources, CI/CD, API specifications, and policy evidence are suite-authored contemporary applications, not assertions that these tools appear in the 2011 book.

\small
\begin{tabularx}{\linewidth}{@{}YYY@{}}
\toprule
\textbf{Artifact} & \textbf{Question answered} & \textbf{Required distinction} \\
\midrule
Architecture description (AD) & What is the system architecture, expressed for stakeholders? & A collection of relevant views plus cross-view information. \\
Architecture view & What does a particular architecture structure look like for this system? & Conforms to a chosen viewpoint and names its element/relation types. \\
Architecture viewpoint specification & How must a particular kind of view be constructed and interpreted? & Reusable conventions and review criteria; not a system-specific view. \\
Architecture decision record (ADR) & Why was a consequential choice made? & A decision, context, alternatives, consequences, status, and links. \\
Template vs. blank & What is the required structure vs. the form to fill? & Template teaches; blank contains neutral labels and writable space. \\
\bottomrule
\end{tabularx}
\normalsize
\section{Architecture structures, styles, views, and viewpoints}
The architectural system is not a single hierarchy. A module decomposition addresses the implementation units and their responsibilities; a component-and-connector (C\&C) view describes runtime elements and their interaction paths; an allocation view relates software elements to nonsoftware structures such as deployment nodes or work assignments. A style constrains permitted types and relations. A view applies selected conventions to one system. A viewpoint specifies the conventions used to construct and evaluate a class of views.

\small
\begin{tabularx}{\linewidth}{@{}YYYY@{}}
\toprule
\textbf{Structure} & \textbf{Primary elements} & \textbf{Representative relations} & \textbf{Questions / useful analyses} \\
\midrule
Module & modules, packages, subsystems & is-part-of, uses, generalizes, layered access & change impact; team boundaries; cyclic dependencies \\
C\&C & runtime components, processes, services; connectors & calls, events, streams, protocol interaction & latency; failure propagation; data flow; trust boundaries \\
Allocation & software elements and environmental resources & deployed-on, installed-in, assigned-to & residency; scaling; ownership; operational dependencies \\
\bottomrule
\end{tabularx}
\normalsize
\section{Stakeholder-driven view selection}
Start from a named decision or stakeholder job rather than a request for diagrams. List candidate views, combine only where semantics remain precise, and prioritize by stakeholder utility and delivery sequence. This reflects the three-stage view-selection approach in Chapter 9.

\begin{enumerate}
\item Identify stakeholders and formulate concrete questions they will use the documentation to answer.
\item Associate each question with the minimum structures and supporting documentation necessary.
\item Propose candidate views and record scope, expected depth, notation, and owner.
\item Combine candidate views only where the combined notation preserves explicit relation semantics.
\item Rank by decision urgency, architectural risk, reuse, maintenance cost, and verification feasibility.
\item Revisit scope when features, responsibilities, and constraints change.
\end{enumerate}
\small
\begin{tabularx}{\linewidth}{@{}YYY@{}}
\toprule
\textbf{Stakeholder} & \textbf{Typical concern} & \textbf{Useful evidence} \\
\midrule
Developer/maintainer & Where to change a feature safely? & decomposition, uses relations, interface contracts \\
Security architect & Where are privileged flows and identities? & C\&C with trust boundaries, deployment allocation, control evidence \\
Operations/SRE & How do failures and changes affect service? & deployment, runtime dependency, failover traces \\
Test/quality & What can be validated and where? & quality scenarios, interface semantics, state traces \\
Product/owner & Which constraints and options are binding? & system overview, decisions, requirements mappings \\
\bottomrule
\end{tabularx}
\normalsize
\section{Standard anatomy of an architecture view}
Each view should carry the following five components. The element catalog can include entries for elements/properties, relations/properties, interfaces, and behavior. A diagram is the primary presentation, not the whole view.

\small
\begin{tabularx}{\linewidth}{@{}YYY@{}}
\toprule
\textbf{Part} & \textbf{Required content} & \textbf{Review test} \\
\midrule
Primary presentation & scope, elements, typed relations, legend, abstractions & Can an unfamiliar reviewer interpret every relation? \\
Element catalog & identifiers, purpose, properties, interfaces, behavior & Can a maintainer locate every named element? \\
Context & boundary and neighboring systems relevant to the style & Can a reader tell external vs. internal dependencies? \\
Variability guide & configurable choices, valid configurations, default behavior & Can an operator identify configuration and constraint effects? \\
Rationale & constraints, rejected alternatives, trade-offs, decision references & Can a future architect explain why this structure exists? \\
\bottomrule
\end{tabularx}
\normalsize
\section{Architecture description: information beyond views}
The architecture description contains document control, a reader roadmap, conventions explaining how views are documented, the system overview, mappings between views, crosscutting rationale, and a directory/glossary. Views should not duplicate stable context when they can reference the common section.

\begin{itemize}
\item Controlled identity: document ID, revision, status, authorship, review date, and change policy.
\item Roadmap: use cases and reader paths from questions to views.
\item System overview: capabilities, users, scope, environment, and architecture drivers.
\item Viewpoint registry: IDs, stakeholders, concerns, styles, diagram notation and semantic rules.
\item Correspondence catalog: explicit many-to-many relationships across module, runtime, and allocation elements.
\item Global rationale: cross-cutting constraints, strategic options, and linked ADRs.
\item Directory: identifiers, glossary, acronyms, ownership, and links to source-of-truth evidence.
\end{itemize}
\section{Model identities, relations, and correspondence rules}
Give elements globally stable IDs, type every relation, and record scope and environment. Use many-to-many mappings rather than assuming each module corresponds to exactly one deployment instance. Track a connection as a first-class entity when protocol, authentication, retries, or message semantics matter.

\small
\begin{tabularx}{\linewidth}{@{}YYY@{}}
\toprule
\textbf{Source / target} & \textbf{Example relationship} & \textbf{Invariants} \\
\midrule
MOD-API -> RUN-API & implemented-as & each production runtime unit maps to released implementation or declared third-party dependency \\
RUN-API -> DEP-COMP & deployed-on & every runtime unit has a target environment and runtime configuration \\
REQ-AUTH -> VIEW-RUNTIME & substantiated-by & every high-risk requirement has an evidence-bearing view or open exception \\
DEC-QUEUE -> RUN-EVENT & realized-by & chosen asynchronous communication must be reflected in connector type \\
\bottomrule
\end{tabularx}
\normalsize
\section{Elements, interfaces, and behavioral documentation}
Describe interaction at a level that lets developers verify assumptions. For each interface, capture provided and required capabilities, endpoint/signature, data semantics, authorization, preconditions, postconditions, invariants, error contract, version policy, and backward compatibility. Static structure alone cannot establish behavior; include trace narratives, sequence or activity diagrams, and state/transition models when needed.

\small
\begin{tabularx}{\linewidth}{@{}YYY@{}}
\toprule
\textbf{Evidence} & \textbf{Core attributes} & \textbf{Typical consumer} \\
\midrule
Interface contract & request/response, semantics, errors, authentication, version & client developer / tester \\
Interaction trace & participants, ordered messages, success and failure alternatives & developer / performance engineer \\
State model & states, events, guards, transitions, invariant violations & reliability / security engineer \\
Data contract & ownership, consistency, schema evolution, retention & integration / data engineer \\
\bottomrule
\end{tabularx}
\normalsize
\section{Rationale, decisions, and architectural risk}
Maintain ADRs as independent decision artifacts linked to affected views and requirements. Do not substitute a list of choices for documented alternatives. Decisions should identify which quality attributes improve and which may deteriorate. Keep superseded decisions discoverable and preserve chronology.

\begin{itemize}
\item Record context, forces, constraints, alternatives, selection, trade-offs, consequences, and lifecycle status.
\item Connect to stakeholder concern and quality scenario; specify how the decision will be falsified or re-evaluated.
\item Track an explicit risk register for unknown performance, failure, security, cost, and team capabilities.
\item Use an evidence pointer for benchmark results, penetration tests, proofs of concept, or operational learning.
\end{itemize}
\section{Quality-attribute scenarios and fitness evidence}
Express a quality concern as source, stimulus, environment, affected artifact, expected response, and quantitative response measure. Make the environment explicit; median latency on a quiet staging environment does not establish tail latency under production load.

\small
\begin{tabularx}{\linewidth}{@{}YYY@{}}
\toprule
\textbf{Attribute} & \textbf{Example measurable scenario} & \textbf{Evidence} \\
\midrule
Availability & one zone lost during peak traffic; restore target capacity in 10 minutes & failover exercise, dashboards \\
Performance & at 500 requests/s, p95 endpoint response under 250 ms & load test traces and methodology \\
Security & a denied role attempts read of restricted object; all attempts are rejected and logged & IAM test, audit record \\
Modifiability & introduce a new consumer with no changes to producer API & change experiment and dependency graph \\
Observability & trace a failed request across three runtimes within 3 minutes & trace sample and runbook \\
\bottomrule
\end{tabularx}
\normalsize
\Note{The numerical targets above are illustrative and are not book-prescribed standards. Adopt targets only after requirements negotiation and measured feasibility.}
\section{Architecture review procedure and evidence discipline}
Chapter 11 concerns fitness of architecture documentation for a specified stakeholder purpose; that is distinct from directly proving that the implemented architecture meets its nonfunctional requirements. A documentation review should establish purpose, review subject and versions, tailored question set, independent inspection, consolidated issues, and follow-up disposition.

\begin{enumerate}
\item Declare the review purpose and stakeholders: development, evaluation, conformance, operations, or change planning.
\item Freeze the artifact versions and define what is inside the review boundary.
\item Build a question set suited to both selected viewpoint conventions and stakeholder jobs.
\item Provide a guided orientation and independent review time; log precise finding IDs and evidence.
\item Reconcile findings by severity, owner, due date, disposition, and re-review threshold.
\item Report closure and document residual risk; do not equate low issue count with verified implementation quality.
\end{enumerate}
\section{Traceability as a graph, not a decorative table}
Use ID-stable directed links among concern, requirement, quality scenario, viewpoint, view, element, connector, decision, test, infrastructure asset, and review finding. Every link must have a defined relation type and independent source. Review both forward and backward coverage. Examples: concerned-by, addressed-in, realized-by, allocated-to, decided-in, verified-by, invalidated-by.

\small
\begin{tabularx}{\linewidth}{@{}YYY@{}}
\toprule
\textbf{Check} & \textbf{Failure condition} & \textbf{Remediation} \\
\midrule
Orphan concern & no view or accepted exclusion maps to a major concern & select/revise viewpoint \\
Unrealized decision & accepted ADR has no implementation evidence link & record realization, status, or exception \\
Unverified scenario & quantified quality criterion has no test or analysis & design a measurement \\
Stale allocation & IaC resource removed but diagram unchanged & regenerate evidence and review mapping \\
Broken correspondence & runtime unit exists with no software or external ownership mapping & classify and document third-party runtime \\
\bottomrule
\end{tabularx}
\normalsize
\section{Architecture documentation as code}
A modern implementation can store prose and diagrams in version control and generate evidence through CI. This is an extension of the book, not part of its original tooling advice. Capture diagrams as code where possible; bind IDs to implementation metadata; lint links and IDs; gate releases on broken cross-view consistency; and publish immutable revisioned outputs.

\begin{itemize}
\item Use ownership and approval rules for viewpoints, architecture descriptions, and ADRs.
\item Keep generated evidence provenance: commit hash, environment, run timestamp, and extractor version.
\item Do not assume CI/IaC manifests fully describe actual runtime behavior; reconcile against observation and source.
\item Protect secrets: omit credentials, access tokens, sensitive topology details, and production data from documents.
\item Create a human-readable changelog for architectural differences, not just regenerated diagrams.
\end{itemize}
\section{Worked example: reusable secure cloud environment}
Scenario: a reusable application platform deploys multiple services into isolated AWS environments using Terraform. Model a provisioning module catalog and dependency edges (module view), request/event and identity flows (C\&C), and placement into networking, compute, data, and operational accounts (allocation). This deliberately avoids presenting an infrastructure drawing as a substitute for all three.

\small
\begin{tabularx}{\linewidth}{@{}YYY@{}}
\toprule
\textbf{View} & \textbf{Small concrete element set} & \textbf{Critical relations} \\
\midrule
Module / decomposition & policy module, network module, workload module, observability module & includes, uses \\
C\&C / runtime & edge endpoint, application runtime, identity service, datastore, telemetry sink & HTTPS calls, token verification, metrics/events \\
Allocation / deployment & account boundary, subnet, execution platform, data service & deployed-on, stores-in, monitored-by \\
\bottomrule
\end{tabularx}
\normalsize
\section{Worked example: change impact and drift}
Suppose the service switches from direct synchronous calls to asynchronous events. An ADR selects the change. Update C\&C connector semantics and failure traces, re-evaluate latency and eventual-consistency scenarios, update ownership contracts and schemas, and remap runtime queues and permissions to allocation artifacts. A complete update also revisits module uses relations and test coverage.

\section{Minimal acceptance criteria for a publishable description}
\begin{itemize}
\item Stakeholders and concerns are named and mapped to at least one reasoned view selection.
\item Each included view conforms to its declared viewpoint and defines element and relation semantics.
\item Every diagram has identifiers, a legend, a boundary, and supporting details or a direct reference.
\item A cross-view mapping and at least one consistency rule are testable.
\item Architecturally significant decisions, exceptions, and risks have stable identifiers.
\item Priority quality scenarios are measurable and connected to evidence.
\item Review status, unresolved findings, revision and owner information are present.
\item Source-of-truth and generated-evidence limitations are plainly stated.
\end{itemize}
\section{Reference crosswalk and further learning}
\small
\begin{tabularx}{\linewidth}{@{}YYY@{}}
\toprule
\textbf{This suite} & \textbf{Primary book foundation} & \textbf{Context} \\
\midrule
Technical reference; viewpoint specifications & Prologue, Chapters 1--5 and 9--10 & types, styles, selections, documentation package \\
Interface and behavioral contracts & Chapters 7--8 & structure plus semantics and dynamic behavior \\
Decision and correspondence workbooks & Chapters 6, 10 & rationale and mappings \\
Review scorecards & Chapter 11 & stakeholder-specific documentation fitness review \\
Notation selection exercises & Appendices A--C & UML, SysML and AADL as candidate representations \\
\bottomrule
\end{tabularx}
\normalsize
Primary source: Paul Clements, Felix Bachmann, Len Bass, David Garlan, James Ivers, Reed Little, Paulo Merson, Robert Nord, and Judith Stafford, Documenting Software Architectures: Views and Beyond, Second Edition, Addison-Wesley, 2011. Its discussion of ISO/IEC 42010 reflects the standard context of its publication; this suite does not assert compliance with any later revision.



\end{document}
